\abstractind
Pendekatan pemodelan lawan dalam permainan berulang umumnya diformulasikan sebagai prediksi satu-langkah terhadap aksi lawan, meskipun interaksi strategis berlangsung dalam horizon yang panjang dan bersifat sekuensial. Kondisi ini berpotensi menimbulkan ketidaksesuaian antara objektif pembelajaran model lawan dan kebutuhan perencanaan multi-langkah agen.

Penelitian ini mengkaji peran horizon prediksi dalam konteks \textit{Iterated Prisoner's Dilemma} dengan mengimplementasikan modul pemodelan lawan berbasis \textit{h-step forecasting} yang dilatih menggunakan kesalahan prediksi distribusi aksi teramati. Analisis difokuskan pada akurasi prediksi multi-langkah serta implikasinya terhadap interaksi strategis jangka pendek.

Hasil eksperimen diharapkan menunjukkan bahwa perluasan horizon prediksi dapat memberikan representasi yang lebih kaya terhadap pola strategi lawan yang bergantung pada trajectory interaksi. Diharapkan ini menyoroti pentingnya mempertimbangkan keselarasan antara objektif prediksi dan horizon perencanaan dalam perancangan sistem pemodelan lawan untuk permainan berulang.
\endabstractind